% Part: many-valued-logic
% Chapter: infinite-valued-logics
% Section: lukasiewicz

\documentclass[../../../include/open-logic-section]{subfiles}

\begin{document}

\olfileid{mvl}{inf}{luk}

\olsection{منطق \L ukasiewicz}

\begin{editorial}
  این بخش پیش‌نویسی بسیار کوتاه دربارهٔ منطق بی‌نهایت‌ارزشی
  \L ukasiewicz است.
\end{editorial}

\begin{defn}\ollabel{def:lukasiewicz} منطق بی‌نهایت‌ارزشی
\L ukasiewicz، یعنی~$\LogLuk[\infty]$، با ماتریس زیر تعریف می‌شود:
\begin{enumerate}
  \item زبان گزاره‌ای استاندارد~$\Lang L_0$ با~$\lnot$، $\land$،
  $\lor$ و~$\lif$.
  \item مجموعهٔ ارزش‌های صدق~$V_\infty$.
  \item $1$ تنها ارزش ممتاز است؛ یعنی~$V^+ = \{1\}$.
  \item توابع صدق با توابع زیر داده می‌شوند:
  \begin{align*}
    \tf{\lnot}[\LogLuk](x) & = 1 - x\\
    \tf{\land}[\LogLuk](x,y) & = \min(x,y)\\
    \tf{\lor}[\LogLuk](x,y) & = \max(x,y)\\
    \tf{\lif}[\LogLuk](x,y) & = \min(1,1-(x-y)) = \begin{cases}
      1 & \text{اگر } x \le y\\
      1-(x-y) & \text{در غیر این صورت.}
    \end{cases}
    \end{align*}
\end{enumerate}
منطق~$m$ارزشی \L ukasiewicz نیز به همین شیوه تعریف می‌شود، جز اینکه
$V = V_m$.
\end{defn}

\begin{prop}
  منطق~$\LogLuk[3]$ که با~\olref[thr][luk]{def:lukasiewicz} تعریف شده،
  همان~$\LogLuk[3]$ تعریف‌شده با~\olref{def:lukasiewicz} است.
\end{prop}

\begin{proof}
  این امر با مقایسهٔ جدول‌های صدق رابط‌ها در
  \olref[thr][luk]{def:lukasiewicz} با جدول‌های صدقی که معادلات
  \olref{def:lukasiewicz} تعیین می‌کنند روشن می‌شود:
  \begin{center}
    \begin{tabular}{c|c} 
      $\tf{\lnot}$ & \\ 
      \hline  
      $1$ & $0$ \\ 
      $1/2$ & $1/2$ \\
      $0$ & $1$ 
    \end{tabular}
    \quad
    \begin{tabular}{c|ccc} 
      $\tf{\land}[\LogLuk[3]]$ & $1$ & $1/2$ & $0$ \\ 
      \hline 
      $1$ & $1$ & $1/2$ & $0$ \\ 
      $1/2$ & $1/2$ & $1/2$ & $0$\\ 
      $0$ & $0$ & $0$ & $0$ 
    \end{tabular}
    \\[2ex]
    \begin{tabular}{c|ccc} 
      $\tf{\lor}[\LogLuk[3]]$ & $1$ & $1/2$ & $0$ \\ 
      \hline 
      $1$ & $1$ & $1$ & $1$ \\ 
      $1/2$ & $1$ & $1/2$ & $1/2$ \\
      $0$ & $1$ & $1/2$ & $0$ 
    \end{tabular}
    \quad
    \begin{tabular}{c|ccc} 
      $\tf{\lif}[\LogLuk[3]]$ & $1$ & $1/2$ & $0$ \\ 
      \hline 
      $1$ & $1$ & $1/2$ & $0$ \\ 
      $1/2$ & $1$ & $1$ & $1/2$  \\ 
      $0$ & $1$ & $1$ & $1$ 
    \end{tabular}
  \end{center} 
\end{proof}

\begin{prop}\ollabel{prop:luk-infty-m}
  اگر~$\Gamma \Entails[\LogLuk[\infty]] !B$، آنگاه برای همهٔ
  $m \ge 2$، $\Gamma \Entails[\LogLuk[m]] !B$.
\end{prop}

\begin{proof}
  تمرین.
\end{proof}

\begin{prob}
  \olref[mvl][inf][luk]{prop:luk-infty-m} را ثابت کنید.
\end{prob}

در واقع، عکس این گزاره نیز برقرار است.

منطق بی‌نهایت‌ارزشی \L ukasiewicz رایج‌ترین منطق فازی است. در ادبیات
منطق فازی، رابط شرطی را اغلب به‌صورت~$\lnot !A \lor !B$ تعریف می‌کنند.
حاصل یک منطق قوی کلینیِ بی‌نهایت‌ارزشی خواهد بود.

\begin{prob}
  نشان دهید~$(p \lif q) \lor (q \lif p)$ همان‌گویی
  $\LogLuk[\infty]$ است.
\end{prob}

\end{document}
